\section{Conclusions}
We introduced hierarchical structure into a supervised CBM along two
axes---part-grounded, coarse-to-fine concepts in the concept space, and a
WordNet taxonomy with an LCA severity metric plus a post-hoc CRM decision rule
in the label space. Compared with an unconstrained black-box, the concept
bottleneck costs both accuracy and severity, the well-documented price of
interpretability.

The two concept-space changes do not pay off equally. Adding the
coarse-to-fine concept hierarchy on top of part conditioning underperforms
our expectations: it does not improve accuracy or severity over plain part
conditioning (Table~\ref{tab:results}), and on the concept consistency test
it is no better, and slightly worse, at landing on the exact edited colour
(Table~\ref{tab:crown-intervention}). Part conditioning alone is the more
promising direction: at a modest accuracy cost, it makes concepts far more
consistent with the local evidence they name.

CRM lowers mistake severity at test time at virtually no accuracy cost on
most variants, provided the classifier it is layered onto gives it a
reasonably faithful posterior. Future work: hierarchy-aware training losses,
real part segmentation in place of keypoint-centred crops, a broader concept
consistency test covering more attributes and edits, and pairing CRM with a
hierarchically structured classifier head.
